% ch3.tex -- Module 7, Chapter 3: Light as a particle (black bodies and the
% photoelectric effect). Dot points M7.IQ3.01, M7.IQ3.02, M7.IQ3.03.
% Figures: figures/ch3_bbcurves.py, figures/ch3_uvcat.py (matplotlib); the rest
% is TikZ / pgfplots below.
\chapter{Light as a Particle: Black Bodies and the Photoelectric Effect}\label{ch3:chap}
\covers{M7.IQ3.01, M7.IQ3.02, M7.IQ3.03}

\inquiry{What evidence supports the particle model of light, and what does this
evidence mean for the quantum model of light?}

\begin{goals}
  \item sketch black body curves and use Wien's law $\lambda_\mathrm{max} = \dfrac{b}{T}$
  \item explain the ultraviolet catastrophe and how Planck's idea, $E = hf$, solved it
  \item describe the photoelectric results that the wave model cannot explain
  \item use $K_\mathrm{max} = hf - \phi$ with conservation of energy and the photon model
  \item read photoelectric graphs to find $h$, $f_0$ and $\phi$
\end{goals}

%% ==================================================================
\section{Black body radiation and Planck's quanta}\label{ch3:sec:blackbody}
\covers{M7.IQ3.01}

\begin{bigidea}
Every warm object gives off EM radiation. A black body is a perfect emitter. Its
spectrum depends only on its temperature: a hotter black body has a shorter peak
wavelength. Physics before 1900 could not explain it. In 1900 Planck did: he said
energy is given out in packets, $E = hf$.
\end{bigidea}

\begin{keywords}
  \kw{Black body}{an ideal object that absorbs all the radiation that hits it. It is also a perfect emitter.}
  \kw{Peak wavelength ($\lambda_\mathrm{max}$)}{the wavelength where a black body curve is highest.}
  \kw{Ultraviolet catastrophe}{the wrong prediction, from physics before 1900, of infinite intensity at short wavelengths.}
  \kw{Quantised}{can only have separate values, in steps.}
  \kw{Quantum (plural: quanta)}{one packet of energy, of size $E = hf$.}
\end{keywords}

Turn on an electric stove. First you feel heat. Then the hotplate glows dull red,
then orange. All objects give off \term{thermal radiation}: EM radiation produced
because of their temperature. Objects contain charged particles (electrons and
protons). These particles vibrate all the time. Vibrating charges give off EM waves
(\cref{ch1:chap}). In a hotter object they vibrate more strongly. So it gives off
more radiation, and more of it has a short wavelength.

\begin{sidebyside}[0.43]
  \centering
  \begin{tikzpicture}[x=0.9cm, y=0.9cm]
    \fill[hscmuted!45] (-0.2,-0.2) rectangle (3.2,2.4);
    \fill[hscink!88] (0,0) rectangle (3,2.2);
    \fill[white] (-0.21,0.98) rectangle (0.01,1.26);
    \draw[ray] (-1.2,1.96) -- (0,1.12);
    \draw[ray] (0,1.12) -- (1.6,0);
    \draw[ray, draw opacity=0.75] (1.6,0) -- (3,0.98);
    \draw[ray, draw opacity=0.5] (3,0.98) -- (1.257,2.2);
    \draw[ray, draw opacity=0.28] (1.257,2.2) -- (0.55,1.705);
    \draw[->, hscray!80!black, line width=0.9pt, decorate,
      decoration={snake, amplitude=1.3pt, segment length=5pt, post length=2.2mm}]
      (-0.1,1.02) -- (-1.35,0.35);
    \node[diagramlabel, anchor=south, font=\footnotesize] at (-1.0,2.0) {light in};
    \node[diagramlabel, anchor=north, align=center, font=\footnotesize] at (-1.05,0.3) {radiation\\out};
    \node[anchor=south, text=white, font=\footnotesize, align=center]
      at (1.55,0.75) {absorbed a little\\at each bounce};
    \node[diagramlabel, anchor=north, font=\footnotesize] at (1.5,-0.25) {hot walls};
    \draw[hscmuted, line width=0.5pt] (0.02,1.3) -- (0.6,2.75)
      node[diagramlabel, anchor=south, font=\footnotesize] {small hole};
  \end{tikzpicture}
  \captionof{figure}{A hollow box with a tiny hole acts as a black body.}
  \label{ch3:fig:cavity}
\nextside
  Real objects also reflect light, which hides their own radiation. So physicists
  study an ideal \term{black body}. It absorbs all the radiation that falls on it
  and reflects none. It is also a perfect emitter: at each temperature it gives
  off the most radiation possible. Its radiation is called
  \term{black body radiation}.\par\smallskip
  A hollow box with a tiny hole is almost a black body (\cref{ch3:fig:cavity}).
  Light that enters bounces around and is absorbed. When the box is hot, the hole
  glows with black body radiation. Stars are also close to black bodies.
\end{sidebyside}

\subsection{The black body curve and Wien's law}

\Cref{ch3:fig:bbcurves} shows black body curves. They are drawn from Planck's
formula, which matches the measured curves. (You meet Planck later in this
section.) The
\term{intensity} is the radiation energy arriving each second on each square metre.
The curve shows how this intensity is shared out among the wavelengths.

\begin{sidebyside}[0.57]
  \centering
  \includegraphics{ch3_bbcurves.pdf}
  \captionof{figure}{Black body curves at three temperatures.}\label{ch3:fig:bbcurves}
\nextside
  \begin{itemize}
    \item Each curve has one peak, at the \term{peak wavelength} $\lambda_\mathrm{max}$.
    \item It rises steeply on the short-wavelength side (left). It falls slowly,
          with a long tail, on the long-wavelength side (right).
    \item A hotter body is higher at \emph{every} wavelength. The curves never cross.
    \item A hotter body peaks at a shorter wavelength.
    \item The curve depends only on temperature, not on the material.
  \end{itemize}
\end{sidebyside}

\begin{refresher}{Kelvin and the EM spectrum}
Absolute temperature $T$ is in kelvin and starts at absolute zero:
$T(\unit{K}) = T(\unit{\degreeCelsius}) + 273$, so $\qty{27}{\degreeCelsius} =
\qty{300}{K}$. Bands: ultraviolet about \qtyrange{10}{400}{nm}, visible
\qtyrange{400}{700}{nm}, infrared \qty{700}{nm} to \qty{1}{mm}.
\end{refresher}

Look at the visible band in \cref{ch3:fig:bbcurves}. A hot object glows with the
colours that its curve covers in this band.
\begin{itemize}
  \item A stove or lava (about \qtyrange{1000}{1500}{K}) peaks in the infrared. Only
        the short-wavelength edge of its curve reaches the visible band, at the red
        end. So it glows dull red.
  \item A very hot object, like the Sun (about \qty{5800}{K}), gives out all the
        visible colours strongly. So it looks white. Even hotter stars look
        blue-white. So a red star is cooler than a blue star.
  \item A cool object, like a person (about \qty{310}{K}), gives out almost no
        visible light. It does not glow.
\end{itemize}

In 1893 Wilhelm Wien showed that the peak wavelength is \emph{inversely
proportional} to the absolute temperature (the temperature in kelvin). This means:
if $T$ doubles, $\lambda_\mathrm{max}$ halves. Later measurements confirmed his law.

\begin{toolkit}{Units for light: nm, \texorpdfstring{\si{\micro m}}{um} and eV}
\begin{itemize}
  \item $\qty{1}{nm} = \qty{1e-9}{m}$ and $\qty{1}{\micro m} = \qty{1e-6}{m}$.
        So $\qty{532}{nm} = \qty{532e-9}{m} = \qty{5.32e-7}{m}$, and
        $\qty{2.0e-6}{m} = \qty{2.0}{\micro m}$.
  \item The \term{electron volt} (eV) is the energy one electron gains when it moves
        through a potential difference (voltage) of \qty{1}{V}.
        $\qty{1}{eV} = \qty{1.602e-19}{J}$ (data sheet).
        J $\to$ eV: divide by \num{1.602e-19}. eV $\to$ J: multiply.
  \item Check: one quantum (packet) of visible light has about \qty{1.8}{eV}
        (red, \qty{700}{nm}) to \qty{3.1}{eV} (violet, \qty{400}{nm}). Infrared
        quanta have less energy. Ultraviolet quanta have more.
\end{itemize}
\end{toolkit}

\begin{toolkit}{Rearranging a formula (reminder)}
To make another symbol the subject, do the same thing to both sides.
\begin{itemize}
  \item $\lambda_\mathrm{max} = \dfrac{b}{T}$. Multiply both sides by $T$:
        $\lambda_\mathrm{max}T = b$. Divide both sides by $\lambda_\mathrm{max}$:
        $T = \dfrac{b}{\lambda_\mathrm{max}}$. Since $\lambda_\mathrm{max}T = b$ at
        every temperature, $\lambda_1T_1 = \lambda_2T_2$.
  \item $K = \frac12 mv^2$. Multiply both sides by 2: $2K = mv^2$. Divide both sides
        by $m$: $v^2 = \dfrac{2K}{m}$. Take the square root: $v = \sqrt{\dfrac{2K}{m}}$.
  \item Check by putting numbers back in: both sides must still be equal.
\end{itemize}
\end{toolkit}

\begin{formulacard}{Wien's displacement law}
  \cardformula{\lambda_\mathrm{max} = \frac{b}{T}}
  \begin{symbols}
    \sym{$\lambda_\mathrm{max}$}{peak wavelength of the black body curve}{\si{m}}
    \sym{$b$}{Wien's displacement constant, \qty{2.898e-3}{m.K} (data sheet)}{\si{m.K}}
    \sym{$T$}{absolute temperature of the object}{\si{K}}
  \end{symbols}
  \inwords{the hotter a black body, the shorter its peak wavelength.}
  \whentouse{black bodies or near black bodies: stars, lava, hot metal, people.}
  \whennottouse{line spectra from a gas lamp, or reflected light.}
  \rearranged{$T = \dfrac{b}{\lambda_\mathrm{max}}$ \qquad
    $\lambda_\mathrm{max}T = b$, so $\lambda_1 T_1 = \lambda_2 T_2$}
  \tensecond{$T = \qty{5800}{K}$, so $\lambda_\mathrm{max} =
    \dfrac{\num{2.898e-3}}{5800} = \qty{5.0e-7}{m} = \qty{500}{nm}$ (visible).}
  \trap{$T$ must be in kelvin, not \si{\degreeCelsius}. The answer is in metres:
    change it to \si{nm} or \si{\micro m} before you name the band.}
\end{formulacard}

\begin{workedexample}{Why does lava glow red?}\label{ch3:we:lava}
Lava is at about \qty{1150}{\degreeCelsius}. Treat it as a black body. Find the
peak wavelength of its radiation and name the EM band. Explain why lava looks red.
\given{$T = 1150 + 273 = \qty{1423}{K}$ \ (Celsius to kelvin), \
  $b = \qty{2.898e-3}{m.K}$}
\find{the peak wavelength $\lambda_\mathrm{max}$ and its band}
\formula{$\lambda_\mathrm{max} = \dfrac{b}{T}$}
\substitute{$\lambda_\mathrm{max} = \dfrac{\qty{2.898e-3}{m.K}}{\qty{1423}{K}}
  = \qty{2.037e-6}{m} = \qty{2.037}{\micro m}$}
\answer{\final{$\lambda_\mathrm{max} = \qty{2.04e-6}{m}$}, which is
  \textbf{infrared} (longer than \qty{700}{nm}). Lava looks red because the
  short-wavelength side of its curve reaches the red end of the visible band. We see
  that small red part. Most of the energy is infrared, which we feel as heat.}
\checkstep{lava is cooler than the Sun (\qty{5800}{K}, peak \qty{500}{nm}), so its
  peak is longer. $5800 \div 1423 \approx 4$, and $\qty{2.04}{\micro m} \approx
  4 \times \qty{500}{nm}$ \cmark.}
\end{workedexample}

\begin{workedexample}{How hot is Sirius? (HSC style)}\label{ch3:we:sirius}
Sirius is the brightest star in the night sky. Its black body curve peaks at
\qty{292}{nm}. (a) Calculate the surface temperature of Sirius. (b) The peak is in
the ultraviolet, which we cannot see. Explain why Sirius still looks bright and
blue-white. (3 marks)
\given{$\lambda_\mathrm{max} = \qty{292}{nm} = \qty{2.92e-7}{m}$ \ (nm to m), \
  $b = \qty{2.898e-3}{m.K}$}
\find{the surface temperature $T$}
\formula{$\lambda_\mathrm{max} = \dfrac{b}{T}$, \ rearranged: $T = \dfrac{b}{\lambda_\mathrm{max}}$}
\substitute{$T = \dfrac{\qty{2.898e-3}{m.K}}{\qty{2.92e-7}{m}} = \qty{9925}{K}$}
\answer{(a) \final{$T = \qty{9.92e3}{K}$}, about \qty{10000}{K}. (b) A hotter body
  is higher at every wavelength, so the curve of Sirius is high across the whole
  visible band. It is highest at the violet and blue end, nearest the peak. So
  Sirius looks bright and blue-white.}
\checkstep{the peak is shorter than the Sun's \qty{500}{nm}, so Sirius must be
  hotter than \qty{5800}{K} \cmark. $500 \div 292 \approx 1.7$ and
  $9925 \div 5800 \approx 1.7$ \cmark. The unit: $\unit{m.K} \div \unit{m} =
  \unit{K}$ \cmark.}
\end{workedexample}

\subsection{The ultraviolet catastrophe}

\begin{sidebyside}[0.56]
  \centering
  \includegraphics{ch3_uvcat.pdf}
  \captionof{figure}{Classical prediction and measured curve at \qty{5000}{K}.}
  \label{ch3:fig:uvcat}
\nextside
  Physicists tried to explain the measured curve with \term{classical physics}:
  physics before 1900, in which light is a wave and energy can have any value.\par\smallskip
  The classical prediction is called the \term{Rayleigh--Jeans law} (after two
  British physicists). It comes closer and closer to the measured curve at long
  wavelengths (\cref{ch3:fig:uvcat}). At short wavelengths it rises without limit.
  It predicts infinite intensity in the ultraviolet. Then the total energy would also
  be infinite, which is impossible. Real curves fall to zero.
\end{sidebyside}

This failure was later called the \term{ultraviolet catastrophe} (a catastrophe is a
disaster). Why did the theory fail? Inside a hot box, the radiation forms standing
waves (waves that fit exactly between the walls). There are more and more possible
standing-wave patterns at shorter wavelengths. Classical physics gives each pattern
the same share of energy. Endless shares add up to infinite energy.

\subsection{Planck's idea: energy comes in packets}

A \term{continuous} quantity can have any value, like the amount of water in a
bottle. A \term{quantised} (or discrete) quantity can only have separate values, in
steps, like the number of eggs in a carton (never 2.5 eggs).

In 1900 Max Planck found a formula that fits the measured curves at every
wavelength. To get it, he made a bold guess about the vibrating charges in the
walls. He called them \term{oscillators}: things that vibrate back and forth. An
oscillator that vibrates with frequency $f$ gives out radiation with the same
frequency $f$. Planck said that such an oscillator can only have the energies
$0, hf, 2hf, 3hf, \ldots$. So it gives out or takes in energy only in packets of
size $hf$. Each packet is a \term{quantum}. The number $h$ is \term{Planck's
constant}, \qty{6.626e-34}{J.s}. It links the size of a packet to its frequency.

\begin{figure}
  \centering
  \begin{tikzpicture}[x=0.9cm, y=0.62cm]
    \begin{scope}
      \draw[axisline] (0,0) -- (0,2.5) node[above, diagramlabel, font=\footnotesize] {energy};
      \fill[hscbrand!14] (0,0) -- (3,0) -- (3,2.0) -- cycle;
      \draw[hscink, line width=0.8pt] (0,0) -- (3,2.0) -- (3,0);
      \fill[object] (1.48,1.12) circle (3.2pt);
      \node[diagramlabel, align=center, anchor=north, font=\footnotesize] at (1.5,-0.1)
        {(a) continuous (ramp):\\any energy};
    \end{scope}
    \begin{scope}[xshift=4.9cm]
      \draw[axisline] (0,0) -- (0,2.5) node[above, diagramlabel, font=\footnotesize] {energy};
      \fill[hscbrand!14] (0,0)
        \foreach \i in {0,...,7} { -- ({0.375*\i},{0.25*(\i+1)}) -- ({0.375*(\i+1)},{0.25*(\i+1)}) }
        -- (3,0) -- cycle;
      \draw[hscink, line width=0.8pt] (0,0)
        \foreach \i in {0,...,7} { -- ({0.375*\i},{0.25*(\i+1)}) -- ({0.375*(\i+1)},{0.25*(\i+1)}) };
      \draw[guide, hscforce, line width=0.8pt] (0,1.2) -- (3.1,1.2);
      \node[anchor=west, text=hscforce, font=\footnotesize, align=left]
        at (3.15,1.2) {heat\\energy\\available};
      \fill[object] (1.31,1.15) circle (3.2pt);
      \draw[hscink, line width=0.4pt] (0.78,0.24) -- (0.43,0.37);
      \node[anchor=west, font=\scriptsize, inner sep=1pt] at (0.78,0.2) {one step $= hf$};
      \node[diagramlabel, align=center, anchor=north, font=\footnotesize] at (1.5,-0.1)
        {(b) low $f$: small steps of $hf$;\\many can be reached};
    \end{scope}
    \begin{scope}[xshift=9.8cm]
      \draw[axisline] (0,0) -- (0,2.5) node[above, diagramlabel, font=\footnotesize] {energy};
      \fill[hscbrand!14] (0,0) -- (1.5,0) -- (1.5,1.7) -- (3,1.7) -- (3,0) -- cycle;
      \draw[hscink, line width=0.8pt] (0,0) -- (1.5,0) -- (1.5,1.7) -- (3,1.7) -- (3,0);
      \draw[guide, hscforce, line width=0.8pt] (0,1.2) -- (3.1,1.2);
      \node[anchor=west, text=hscforce, font=\footnotesize, align=left]
        at (3.15,1.2) {heat\\energy\\available};
      \draw[<->, hscink, line width=0.5pt] (1.64,0.04) -- (1.64,1.66);
      \node[anchor=west, font=\scriptsize, align=left, inner sep=1pt] at (1.7,0.45)
        {one step\\$= hf$};
      \fill[object] (0.75,0.14) circle (3.2pt);
      \node[diagramlabel, align=center, anchor=north, font=\footnotesize] at (1.5,-0.1)
        {(c) high $f$: one step of $hf$\\is too big, rarely reached};
    \end{scope}
  \end{tikzpicture}
  \caption{Continuous energy (a) and quantised energy (b, c). The height of the ball
    is the energy of one oscillator. The dashed line is the typical energy that heat
    gives one oscillator.}\label{ch3:fig:stairs}
\end{figure}

How do packets fix the curve (\cref{ch3:fig:stairs})?
\begin{itemize}
  \item Low frequency (long wavelength): each step $hf$ is tiny, so the steps act
        like a ramp. Planck's result agrees with classical physics there.
  \item High frequency (short wavelength): each step $hf$ is large. Heat rarely gives
        one oscillator enough energy for even one step. So little short-wavelength
        radiation is given off, and the curve falls to zero.
  \item A hotter object has more energy to share, so it reaches bigger steps. Its
        peak moves to shorter wavelengths (Wien's law).
\end{itemize}

Planck saw his packets as a maths trick. He still thought light was a wave. But
his idea started quantum physics: light energy is \emph{given out} in packets of
$hf$. Einstein took the next step in 1905. He said that light itself travels as
these packets. One packet (quantum) of light is called a \term{photon}
(\cref{ch3:sec:photon}).

\begin{formulacard}{Energy of one quantum (photon)}
  \cardformula{E = hf}
  \begin{symbols}
    \sym{$E$}{energy of one quantum (one photon)}{\si{J}}
    \sym{$h$}{Planck's constant, \qty{6.626e-34}{J.s} (data sheet)}{\si{J.s}}
    \sym{$f$}{frequency of the radiation}{\si{Hz}}
    \sym{$\lambda$}{wavelength of the radiation (in the forms with $c$)}{\si{m}}
    \sym{$c$}{speed of light in a vacuum, \qty{3.00e8}{m.s^{-1}} (data sheet)}{\si{m.s^{-1}}}
  \end{symbols}
  \inwords{the energy of one packet depends only on its frequency. Higher frequency
    means a bigger packet.}
  \whentouse{the energy of one quantum or one photon of any EM radiation.}
  \whennottouse{the total energy of a beam. That is (number of photons) $\times\, hf$.}
  \rearranged{$f = \dfrac{E}{h}$ \qquad $E = \dfrac{hc}{\lambda}$ \qquad
    $\lambda = \dfrac{hc}{E}$}
  \tensecond{$f = \qty{5.0e14}{Hz}$, so $E = \num{6.626e-34} \times \num{5.0e14}
    = \qty{3.3e-19}{J}$.}
  \derive{$E = \dfrac{hc}{\lambda}$ is \emph{not} printed on the sheet. Combine two
    sheet formulas:
    \begin{align*}
      c &= f\lambda && \text{$v = f\lambda$ with $v = c$ for light}\\
      f &= \frac{c}{\lambda} && \text{divide both sides by $\lambda$}\\
      E &= hf = \frac{hc}{\lambda} && \text{replace $f$ in $E = hf$}
    \end{align*}}
  \trap{change \si{nm} to \si{m} before you use $E = \dfrac{hc}{\lambda}$. The answer
    is in joules. It is in eV only after you divide by \num{1.602e-19}.}
\end{formulacard}



\begin{workedexample}{Why so little ultraviolet? (HSC style)}\label{ch3:we:quanta}
An oven at \qty{1500}{K} gives off radiation that peaks near \qty{2.00}{\micro m}.
Calculate the energy, in J and eV, of one quantum at \qty{2.00}{\micro m} and of one
at \qty{200}{nm}. Use your answers and Planck's idea to explain why the oven gives off
almost no ultraviolet. (4 marks)
\given{$\lambda_1 = \qty{2.00}{\micro m} = \qty{2.00e-6}{m}$, \
  $\lambda_2 = \qty{200}{nm} = \qty{2.00e-7}{m}$,
  \ $h = \qty{6.626e-34}{J.s}$, \ $c = \qty{3.00e8}{m.s^{-1}}$}
\find{the two quantum energies, and how they compare}
\formula{$E = \dfrac{hc}{\lambda}$, \ then divide by \num{1.602e-19} for eV}
\substitute{\eqlines{E_1 &= \frac{\qty{6.626e-34}{J.s} \times \qty{3.00e8}{m.s^{-1}}}{\qty{2.00e-6}{m}}
    = \qty{9.94e-20}{J}\\
  E_2 &= \frac{\qty{6.626e-34}{J.s} \times \qty{3.00e8}{m.s^{-1}}}{\qty{2.00e-7}{m}}
    = \qty{9.94e-19}{J}}\par
  \eqlines{E_1 &= \frac{\num{9.94e-20}}{\num{1.602e-19}} = \qty{0.620}{eV} \quad\text{(J to eV)}\\
  E_2 &= \frac{\num{9.94e-19}}{\num{1.602e-19}} = \qty{6.20}{eV} \quad\text{(J to eV)}\\
  \frac{E_2}{E_1} &= \frac{\num{9.94e-19}}{\num{9.94e-20}} = 10}}
\answer{\final{$E_1 = \qty{9.94e-20}{J} = \qty{0.620}{eV}$}, \
  \final{$E_2 = \qty{9.94e-19}{J} = \qty{6.20}{eV}$}. A UV quantum is 10 times
  bigger. An oscillator must give out a whole quantum at once. Heat sometimes gives
  one oscillator enough energy for one infrared quantum. It almost never gives one
  oscillator 10 times as much energy at once. So almost no UV quanta are given out.}
\checkstep{$E \propto \dfrac{1}{\lambda}$ ($\propto$ means ``is proportional
  to''). $\lambda$ is 10 times smaller, so $E$ is 10 times bigger \cmark. Wien:
  $\dfrac{b}{1500} = \qty{1.93}{\micro m}$, close to \qty{2.00}{\micro m} \cmark.}
\end{workedexample}

\begin{mistakes}
\begin{itemize}
  \item \xmark\ Putting \si{\degreeCelsius} into $\lambda_\mathrm{max} = \dfrac{b}{T}$.
        \cmark\ Add 273 first.
  \item \xmark\ ``A cooler object gives out more red light than a hotter one.''
        \cmark\ The hotter curve is higher at every wavelength. Only the cooler
        object's \emph{peak} is redder.
  \item \xmark\ Drawing a symmetric bell curve, or curves that cross.
        \cmark\ A steep rise on the short-wavelength side, a long tail on the
        long-wavelength side. The hotter curve is above everywhere.
  \item \xmark\ ``Planck said light is made of photons.'' \cmark\ Planck quantised
        the oscillators in the walls. Einstein quantised light itself.
\end{itemize}
\end{mistakes}

\begin{examtip}
``Explain how Planck's hypothesis solved the ultraviolet catastrophe'' (4 marks).
One mark for each linked point: (1) classical theory predicted intensity rising
without limit at short wavelengths, but measured curves peak and fall to zero;
(2) Planck assumed oscillators give out energy only in quanta, $E = hf$; (3) at high
$f$ the quanta are large compared with the heat energy available, so few are given
out; (4) his formula fitted the measured curves at all wavelengths, and began the
quantum model of light. In sketches, put the hotter curve higher everywhere, with
its peak to the left.
\end{examtip}

\begin{quickcheck}
  \item A student says: ``A black body is any object that is painted black.''
        Explain why this is wrong, and give the correct meaning of a black body.
        (2 marks)\label{ch3:qc:bbdef}
  \item An electric heater element is at \qty{600}{\degreeCelsius}. Find the peak
        wavelength of its radiation, and name the EM band.
  \item The radiation from a lamp filament peaks at \qty{1.05}{\micro m}. Find the
        temperature of the filament in kelvin and in \si{\degreeCelsius}.
  \item Calculate the energy of a photon of violet light (\qty{400}{nm}) in joules
        and in electron volts.
  \item For which part of the black body curve did classical physics agree with
        experiment? Why was its failure called the ``ultraviolet''
        catastrophe?\label{ch3:qc:uvname}
\end{quickcheck}

%% ==================================================================
\section{The photoelectric effect: evidence against the wave model}\label{ch3:sec:photoelectric}
\covers{M7.IQ3.02}

\begin{bigidea}
Light that shines on a metal can knock electrons out of it. The wave model of light
predicts how this should depend on brightness and colour. Experiments showed that
three of these predictions are wrong.
\end{bigidea}

\begin{keywords}
  \kw{Photoelectric effect}{electrons are released from a metal surface when light shines on it.}
  \kw{Photocurrent}{the current made by the released electrons (photoelectrons).}
  \kw{Threshold frequency ($f_0$)}{the lowest frequency of light that releases electrons from a metal.}
  \kw{Stopping voltage ($V_0$)}{the smallest reverse voltage that stops every photoelectron, so the current is zero.}
\end{keywords}

Shine ultraviolet light on clean zinc, and electrons fly out of the zinc. This is
the \term{photoelectric effect}. The released electrons are \term{photoelectrons}.

Heinrich Hertz first noticed the effect in 1887. He was making Maxwell's radio waves
with sparks. A spark jumps across a small air gap between two metal balls, called a
\term{spark gap}. Hertz saw that the sparks jumped more easily when ultraviolet light
shone on the gap. In 1902 Philipp Lenard measured the effect carefully with a vacuum
tube (\cref{ch3:fig:apparatus}).

\begin{figure}
  \centering
  \begin{tikzpicture}[x=1cm, y=1cm]
    \draw[hscink, line width=0.9pt, fill=hscbrand!5, rounded corners=14pt]
      (0,0) rectangle (6,2.6);
    \draw[hscvel!70, line width=2.6pt] (1.15,2.6) -- (2.75,2.6);
    \node[diagramlabel, anchor=south east, align=right, font=\footnotesize] at (1.1,2.65)
      {quartz window\\(lets UV through)};
    \draw[hscink, line width=2.4pt] (0.8,0.45) -- (0.8,2.15);
    \draw[hscink, line width=2.4pt] (5.2,0.6) -- (5.2,2.0);
    \foreach \ye/\xs in {1.9/2.42, 1.45/2.78, 1.0/3.14} {
      \draw[->, hscray, line width=1pt, decorate,
        decoration={snake, amplitude=1.3pt, segment length=5pt, post length=2.2mm}]
        (\xs,3.8) -- (0.92,\ye);
    }
    \node[diagramlabel, anchor=west, font=\footnotesize] at (3.2,3.75) {light: frequency $f$};
    \foreach \p in {(1.9,1.65), (3.0,1.0), (4.0,1.75)} {
      \fill[hscvel] \p circle (2.6pt);
      \draw[velocity, line width=0.9pt] \p ++(0.12,0) -- ++(0.55,0);
    }
    \node[text=hscvel, font=\footnotesize] at (3.0,0.62) {photoelectrons};
    \node[diagramlabel, anchor=north, font=\footnotesize] at (3.6,2.45) {vacuum};
    \draw[hscmuted, line width=0.5pt] (0.75,1.6) -- (-0.2,1.6)
      node[diagramlabel, anchor=east, align=right, font=\footnotesize] {metal\\cathode};
    \draw[hscmuted, line width=0.5pt] (5.25,1.6) -- (6.2,1.6)
      node[diagramlabel, anchor=west, align=left, font=\footnotesize] {anode\\(collector)};
    % circuit
    \draw[hscink, line width=0.8pt] (0.8,0.45) -- (0.8,-1.0) -- (2.2,-1.0);
    \draw[hscink, line width=0.8pt] (2.4,-1.0) -- (3.67,-1.0);
    \draw[hscink, line width=0.8pt] (4.33,-1.0) -- (5.2,-1.0) -- (5.2,0.6);
    \draw[hscink, line width=0.9pt] (2.2,-1.35) -- (2.2,-0.65);
    \draw[hscink, line width=2.2pt] (2.4,-1.18) -- (2.4,-0.82);
    \draw[->, hscink, line width=0.7pt] (1.9,-1.4) -- (2.7,-0.55);
    \node[font=\scriptsize] at (2.02,-0.52) {$+$};
    \node[font=\scriptsize] at (2.6,-1.3) {$-$};
    \draw[hscink, line width=0.8pt] (1.6,-1.0) -- (1.6,-1.8) -- (2.0,-1.8);
    \draw[hscink, line width=0.8pt] (2.6,-1.8) -- (3.0,-1.8) -- (3.0,-1.0);
    \draw[hscink, line width=0.8pt, fill=white] (2.3,-1.8) circle (0.3) node[font=\small] {V};
    \node[diagramlabel, anchor=north, font=\footnotesize, align=center] at (2.3,-2.15)
      {variable supply (can be reversed)\\with a voltmeter};
    \draw[hscink, line width=0.8pt, fill=white] (4.0,-1.0) circle (0.33)
      node[font=\small] {\si{\micro A}};
    \node[diagramlabel, anchor=south, font=\footnotesize] at (4.0,-0.64) {microammeter};
  \end{tikzpicture}
  \caption{Photoelectric apparatus. Light releases electrons from the metal cathode.
    They cross the vacuum to the anode. Here the supply makes the anode negative, as
    used to find $V_0$. Reverse the supply to make the anode positive and collect
    every electron.}\label{ch3:fig:apparatus}
\end{figure}

\begin{itemize}
  \item Light falls on the \term{cathode}, the metal plate that releases electrons.
        They cross to the \term{anode}, the collecting plate. The whole tube is
        often called a \term{photocell}.
  \item The \term{microammeter} (a meter for very small currents, in millionths of
        an ampere) measures the \term{photocurrent}. A bigger current means more
        electrons each second, not faster electrons.
  \item The vacuum stops electrons from hitting air molecules. The quartz window lets
        ultraviolet in (ordinary glass absorbs it).
\end{itemize}

\begin{sidebyside}[0.4]
  \centering
  \begin{tikzpicture}[x=1.1cm, y=1.1cm]
    % glass case and stand
    \draw[hscink, line width=0.7pt, fill=hscbrand!4, rounded corners=6pt]
      (-0.8,-1.6) rectangle (0.8,0.25);
    \fill[hscmuted!50] (-1.0,-1.85) rectangle (1.0,-1.6);
    % rod and zinc plate
    \draw[hscink, line width=1.6pt] (0,-1.35) -- (0,0.75);
    \fill[hscmuted!70] (-0.8,0.75) rectangle (0.8,0.9);
    \draw[hscink, line width=0.6pt] (-0.8,0.75) rectangle (0.8,0.9);
    \foreach \x in {-0.6,-0.2,0.2,0.6}
      \node[font=\scriptsize, text=hscvel, inner sep=0pt] at (\x,1.0) {$-$};
    \draw[hscmuted, line width=0.5pt] (0.8,0.82) -- (1.15,0.82)
      node[diagramlabel, anchor=west, font=\footnotesize] {zinc plate};
    % leaf standing out (charged) and fallen (charge gone)
    \draw[hscacc, line width=1.4pt] (0,-0.3) -- (0.5,-1.0);
    \draw[hscacc, line width=1.1pt, dashed] (0.06,-0.3) -- (0.1,-1.15);
    \draw[hscmuted, line width=0.5pt] (0.35,-0.8) -- (1.15,-0.2)
      node[diagramlabel, anchor=west, font=\footnotesize, align=left]
      {leaf out\\(charged)};
    \draw[hscmuted, line width=0.5pt] (0.12,-1.1) -- (1.15,-1.35)
      node[diagramlabel, anchor=west, font=\footnotesize, align=left]
      {leaf fallen\\(charge gone)};
    \node[diagramlabel, anchor=north, font=\footnotesize] at (0,-1.88) {electroscope};
    % UV lamp
    \fill[hscfield!60] (-1.95,1.95) rectangle (-1.25,2.3);
    \node[diagramlabel, anchor=south, font=\footnotesize] at (-1.6,2.3) {UV lamp};
    \foreach \i in {0,1,2}
      \draw[->, hscfield, line width=0.9pt, decorate,
        decoration={snake, amplitude=1.1pt, segment length=4pt, post length=1.8mm}]
        ({-1.4+0.2*\i},1.9) -- ({-0.7+0.42*\i},1.2);
  \end{tikzpicture}
  \captionof{figure}{A gold-leaf electroscope with a zinc plate. UV light makes the
    leaf fall.}\label{ch3:fig:electroscope}
\nextside
  A classroom demonstration uses a gold-leaf \term{electroscope}
  (\cref{ch3:fig:electroscope}). It is a metal rod with a thin gold leaf hanging from
  it, inside a glass case. When the electroscope is charged, the leaf is pushed away
  from the rod and stands out. When the charge leaves, the leaf falls back.\par\smallskip
  Put a clean zinc plate on top and give it a negative charge (extra electrons).
  Shine ultraviolet light on the zinc. The leaf falls, because electrons leave the
  zinc. A very bright visible lamp does nothing. A sheet of glass, which absorbs UV,
  stops the effect. So \emph{whether} electrons leave depends on the frequency of the
  light, not on its brightness.
\end{sidebyside}

\subsection{Stopping voltage}

Now make the anode negative, as drawn in \cref{ch3:fig:apparatus}. The negative anode repels the
electrons: it pushes them back. As the reverse voltage increases, fewer electrons
reach the anode and the current falls. At the \term{stopping voltage} $V_0$, even the
fastest electrons stop just before the anode, and the current is zero.

Electrons leave with a range of speeds. Some electrons start deeper inside the
metal. They lose some energy on the way out, so they leave more slowly. The fastest
electrons have the \term{maximum kinetic energy} $K_\mathrm{max}$. The stopping
voltage measures it.

\begin{refresher}{Work, voltage and kinetic energy}
When a charge $q$ moves through a potential difference (voltage) $V$, the electric
field does work $W = qV$ on it. Work is energy transferred, in joules (Module 6).
Kinetic energy is $K = \frac12 mv^2$ (Year 11). Rearranged (see the toolkit
\emph{Rearranging a formula}): $v = \sqrt{\dfrac{2K}{m}}$, with $K$ in joules. For an
electron, use the size of the charge, \qty{1.602e-19}{C}.
\end{refresher}

\begin{formulacard}[combine={$W = qV$}]{Stopping voltage and maximum kinetic energy}
  \cardformula{K_\mathrm{max} = qV_0}
  \begin{symbols}
    \sym{$K_\mathrm{max}$}{maximum kinetic energy of the photoelectrons}{\si{J}}
    \sym{$q$}{size of the electron's charge, \qty{1.602e-19}{C}}{\si{C}}
    \sym{$V_0$}{stopping voltage}{\si{V}}
  \end{symbols}
  \inwords{the stopping voltage takes away all the kinetic energy of the fastest
    electrons.}
  \whentouse{you know $V_0$ and need $K_\mathrm{max}$, or the other way round.}
  \whennottouse{for a typical electron. Most electrons have less than $K_\mathrm{max}$.}
  \rearranged{$V_0 = \dfrac{K_\mathrm{max}}{q}$. \ $K_\mathrm{max}$ in eV is the
    same number as $V_0$ in volts.}
  \tensecond{$V_0 = \qty{2.0}{V}$, so $K_\mathrm{max} = \num{1.602e-19} \times 2.0
    = \qty{3.2e-19}{J} = \qty{2.0}{eV}$.}
  \derive{the reverse voltage does work against each electron. The fastest one just
    stops at the anode, so all its kinetic energy is used up:
    \begin{align*}
      W &= qV_0 && \text{work done against the electron ($W = qV$)}\\
      qV_0 &= K_\mathrm{max} && \text{work done = kinetic energy lost}
    \end{align*}}
  \trap{$V_0$ is a voltage (V). $K_\mathrm{max}$ is an energy (J or eV). Never
    write ``$K_\mathrm{max} = \qty{2.0}{V}$''.}
\end{formulacard}

\begin{workedexample}{From stopping voltage to electron speed}\label{ch3:we:stopping}
In a photoelectric experiment the stopping voltage is \qty{1.40}{V}. Calculate the
maximum kinetic energy of the photoelectrons (in J and eV) and their maximum speed.
\given{$V_0 = \qty{1.40}{V}$, \ $q = \qty{1.602e-19}{C}$, \ $m_\mathrm{e} = \qty{9.109e-31}{kg}$}
\find{$K_\mathrm{max}$ (in J and eV) and $v_\mathrm{max}$}
\formula{$K_\mathrm{max} = qV_0$, \ then $K = \tfrac12 mv^2$ rearranged:
  $v = \sqrt{\dfrac{2K}{m}}$}
\substitute{\eqlines{K_\mathrm{max} &= \qty{1.602e-19}{C} \times \qty{1.40}{V}
    = \qty{2.243e-19}{J} = \qty{1.40}{eV}\\
  v_\mathrm{max} &= \sqrt{\frac{2 \times \qty{2.243e-19}{J}}{\qty{9.109e-31}{kg}}}
    = \qty{7.017e5}{m.s^{-1}}}}
\answer{\final{$K_\mathrm{max} = \qty{2.24e-19}{J} = \qty{1.40}{eV}$}, \
  \final{$v_\mathrm{max} = \qty{7.02e5}{m.s^{-1}}$}}
\checkstep{the eV number equals the volts number \cmark. The speed is below 1\,\%
  of $c$, so $K = \frac12 mv^2$ is fine \cmark.}
\end{workedexample}

\subsection{What the wave model predicted}

In the wave model, the energy of light depends on its brightness (amplitude), not on
its frequency. The energy arrives smoothly, spread over the whole surface.
\Cref{ch3:tab:wave} compares the predictions with what was observed.

\begin{table}
  \caption{Wave-model predictions and photoelectric observations.}\label{ch3:tab:wave}
  \centering
  \small
  \begin{tabular}{@{}L{52mm}L{64mm}L{20mm}@{}}
    \toprule
    \colhead{The wave model predicted} & \colhead{Experiments showed} & \colhead{Wave model}\\
    \midrule
    Brighter light gives more electrons each second, so a bigger current.
      & Brighter light gives a bigger current (when $f \ge f_0$).
      & \cmark\ agrees\\
    Any frequency works if the light is bright enough.
      & There is a threshold frequency $f_0$. Below it, no electrons leave, however
        bright the light.
      & \xmark\ fails\\
    With dim light, electrons need time to build up energy: a delay.
      & No delay. Electrons leave at once (in less than \qty{1e-9}{s}), even in very
        dim light.
      & \xmark\ fails\\
    Brighter light makes faster electrons, so a bigger $V_0$.
      & $V_0$ (and $K_\mathrm{max}$) does not depend on brightness. It increases
        in a straight line with frequency.
      & \xmark\ fails\\
    \bottomrule
  \end{tabular}
\end{table}

\begin{workedexample}{The wave model's waiting time (HSC style)}\label{ch3:we:delay}
Very dim light of intensity \qty{1.0e-2}{W.m^{-2}} falls on sodium. An electron needs
at least \qty{2.3}{eV} to escape. In the wave model, the energy is spread evenly over
the surface. Assume one electron collects only the energy that falls on one atom, an
area of about \qty{1.0e-19}{m^2}. (a) Estimate how long the wave model says the
electron must wait. (b) Real electrons leave in less than \qty{1e-9}{s}. What does
this tell us?
\given{intensity $= \qty{1.0e-2}{W.m^{-2}}$ (energy each second on each \si{m^2}),
  \ $A = \qty{1.0e-19}{m^2}$, \ $\Delta E = \qty{2.3}{eV}$}
\find{the waiting time $\Delta t$}
\formula{power on one atom: $P = \text{intensity} \times A$.\par
  $P = \dfrac{\Delta E}{\Delta t}$, where $\Delta$ means ``change in''.\par
  Multiply both sides by $\Delta t$: $P\,\Delta t = \Delta E$.\par
  Divide both sides by $P$: $\Delta t = \dfrac{\Delta E}{P}$}
\substitute{\eqlines{P &= \qty{1.0e-2}{W.m^{-2}} \times \qty{1.0e-19}{m^2}
    = \qty{1.0e-21}{W}\\
  \Delta E &= 2.3 \times \qty{1.602e-19}{J} = \qty{3.685e-19}{J}
    \quad\text{(eV to J)}\\
  \Delta t &= \frac{\qty{3.685e-19}{J}}{\qty{1.0e-21}{W}} = \qty{368.5}{s}}}
\answer{(a) \final{$\Delta t \approx \qty{3.7e2}{s}$}, about 6 minutes. (b) There is no
  delay, so the energy cannot arrive spread out like a wave. It must arrive in
  concentrated packets. This is evidence against the wave model.}
\checkstep{$\unit{J} \div \unit{W} = \unit{J} \div \unit{J.s^{-1}} = \unit{s}$ \cmark.
  Dimmer light would make the wave-model wait even longer.}
\end{workedexample}

\begin{mistakes}
\begin{itemize}
  \item \xmark\ ``Photons are emitted from the metal.''
        \cmark\ Photons arrive. \emph{Photoelectrons} leave.
  \item \xmark\ ``Brighter light makes faster electrons.''
        \cmark\ Brighter light makes more electrons, not faster ones.
  \item \xmark\ ``The wave model got everything wrong.''
        \cmark\ It correctly predicted that brighter light gives a bigger current.
  \item \xmark\ Using $q = \qty{-1.602e-19}{C}$ and getting a negative energy.
        \cmark\ Use the size of the charge.
\end{itemize}
\end{mistakes}

\begin{examtip}
``Explain why photoelectric results did not support the wave model'' (3--5 marks).
For each observation: (1) state it (for example, a threshold frequency); (2) state
what the wave model predicted (any frequency works if bright enough, because wave
energy depends on intensity); (3) say that they disagree. For 3 marks, give two
observations, each linked to the wave-model prediction it breaks. For 4--5 marks,
give three, and finish with the conclusion: light energy must arrive in packets of
$hf$. A list of observations with no link to the wave model does not earn full
marks.
\end{examtip}

\begin{quickcheck}
  \item The stopping voltage in an experiment is \qty{0.85}{V}. Find the maximum
        kinetic energy of the photoelectrons in eV and in joules.
  \item Why is the photocell in \cref{ch3:fig:apparatus} a vacuum, and why does it
        have a quartz window?
  \item Light of the same frequency is made brighter. What happens to (a) the
        photocurrent and (b) the stopping voltage? Which did the wave model predict
        correctly?
  \item Describe two photoelectric observations that the wave model could not
        explain. For each, say what the wave model predicted.\label{ch3:qc:twoobs}
\end{quickcheck}

%% ==================================================================
\section{Einstein's photon model and photoelectric graphs}\label{ch3:sec:photon}
\covers{M7.IQ3.03, M7.IQ3.02}

\begin{bigidea}
In 1905 Einstein said that light itself is a stream of packets, called photons, each
with energy $hf$. One photon gives all its energy to one electron, so conservation
of energy gives $K_\mathrm{max} = hf - \phi$. A graph of $K_\mathrm{max}$ against $f$
is a straight line with gradient $h$.
\end{bigidea}

\begin{keywords}
  \kw{Photon}{one packet (quantum) of light, with energy $E = hf$.}
  \kw{Work function ($\phi$)}{the smallest energy needed to remove an electron from a metal's surface.}
  \kw{Threshold wavelength ($\lambda_0$)}{the longest wavelength that releases electrons from a metal.}
  \kw{Saturation current}{the largest photocurrent, when every released electron reaches the anode.}
  \kw{Wave--particle duality}{light acts like a wave in some experiments and like particles in others.}
\end{keywords}

Planck said the walls give out energy in packets. Einstein said that light
\emph{itself} is a stream of packets. Each packet is a photon, with energy
$E = hf$. A brighter beam of one colour has \emph{more photons} each second. Each
photon still has the same energy.

Einstein's key rule: \textbf{one photon interacts with one electron.} The photon
gives all of its energy to that electron, at once. An electron cannot collect energy
from two photons.

Think of a drink machine that takes one coin at a time and never adds coins
together. A drink costs \$1.50. A \$2 coin works, and you get 50 cents change.
Fifteen 10-cent coins never work, however many you put in. The coin is like one
photon. The price is like the work function $\phi$ (below). The change is like the
electron's kinetic energy $K_\mathrm{max}$.

To escape, an electron must do work against the pull of the metal. The smallest
energy it needs is the \term{work function} $\phi$. Each metal has its own value: for
example, caesium \qty{2.1}{eV} and zinc \qty{4.3}{eV}.

\begin{figure}
  \centering
  \begin{tikzpicture}[x=1cm, y=1cm]
    \node[diagramlabel, anchor=east] at (-0.15,1.55) {$hf > \phi$:};
    \fill[hscmuted!30] (0,1.2) rectangle (5.46,1.9);
    \fill[hscseriesC!35] (5.46,1.2) rectangle (7.96,1.9);
    \draw[hscink, line width=0.7pt] (0,1.2) rectangle (7.96,1.9);
    \draw[hscink, line width=0.7pt] (5.46,1.2) -- (5.46,1.9);
    \node[font=\footnotesize] at (2.73,1.55) {work function $\phi$ (to escape)};
    \node[font=\footnotesize] at (6.71,1.55) {$K_\mathrm{max}$ (left over)};
    \draw[hscink, line width=0.6pt] (0,2.0) -- (0,2.15) -- (7.96,2.15) -- (7.96,2.0);
    \node[diagramlabel, anchor=south, font=\footnotesize] at (3.98,2.15) {energy of one photon, $hf$};
    \node[diagramlabel, anchor=east] at (-0.15,0.35) {$hf < \phi$:};
    \fill[hscseriesA!25] (0,0) rectangle (4.68,0.7);
    \draw[hscink, line width=0.7pt] (0,0) rectangle (4.68,0.7);
    \draw[hscink, line width=0.7pt, dashed] (4.68,0) -- (5.46,0) -- (5.46,0.7) -- (4.68,0.7);
    \node[font=\footnotesize] at (2.34,0.35) {photon energy $hf$};
    \node[diagramlabel, anchor=west, align=left, font=\footnotesize] at (5.6,0.35)
      {dashed: energy still missing\\to reach $\phi$. No electron leaves.};
  \end{tikzpicture}
  \caption{The energy budget for one photon and one electron.}\label{ch3:fig:budget}
\end{figure}

\begin{refresher}{Conservation of energy (Year 11)}
Energy cannot be made or destroyed. It can only change from one form to another. So
the energy that goes in equals the energy that comes out.
\end{refresher}

Apply conservation of energy to one photon and the electron that is easiest to
remove, the \term{least-bound electron} (\cref{ch3:fig:budget}):
\begin{align}
  \text{energy in} &= \text{energy to escape} + \text{energy left over}
     && \text{conservation of energy}\notag\\
  hf &= \phi + K_\mathrm{max} && \text{put in the symbols}\notag\\
  K_\mathrm{max} &= hf - \phi && \text{subtract $\phi$ from both sides}\label{ch3:eq:einstein}
\end{align}
Electrons deeper in the metal need more than $\phi$, so they leave slower.

The photon model explains every result in \cref{ch3:tab:wave}:
\begin{itemize}
  \item \textbf{Brighter light, bigger current:} more photons each second free more
        electrons each second.
  \item \textbf{Threshold frequency:} one photon must give at least $\phi$, so
        $hf \ge \phi$. Below $f_0$ no single photon is enough, and energies cannot add up.
  \item \textbf{No delay:} each photon delivers its whole energy in one event.
  \item \textbf{$K_\mathrm{max}$ depends on $f$, not brightness:}
        $K_\mathrm{max} = hf - \phi$. Brightness changes the number of photons, not $hf$.
\end{itemize}

\begin{formulacard}{Einstein's photoelectric equation}
  \cardformula{K_\mathrm{max} = hf - \phi}
  \begin{symbols}
    \sym{$K_\mathrm{max}$}{maximum kinetic energy of the photoelectrons}{\si{J} (or eV)}
    \sym{$h$}{Planck's constant, \qty{6.626e-34}{J.s}}{\si{J.s}}
    \sym{$f$}{frequency of the light}{\si{Hz}}
    \sym{$\phi$}{work function of the metal}{\si{J} (or eV)}
    \sym{$\lambda$}{wavelength of the light (in the form with $c$)}{\si{m}}
    \sym{$c$}{speed of light in a vacuum, \qty{3.00e8}{m.s^{-1}}}{\si{m.s^{-1}}}
  \end{symbols}
  \inwords{energy of one photon, minus the energy to escape, equals the kinetic energy
    left over.}
  \whentouse{light (or other EM radiation) falls on a metal, and $hf \ge \phi$.}
  \whennottouse{$hf < \phi$. A negative answer means no electrons at all.}
  \rearranged{$\phi = hf - K_\mathrm{max}$ \qquad $f = \dfrac{K_\mathrm{max} + \phi}{h}$
    \qquad $K_\mathrm{max} = \dfrac{hc}{\lambda} - \phi$}
  \tensecond{photon energy $hf = \qty{3.5}{eV}$ and $\phi = \qty{2.3}{eV}$, so
    $K_\mathrm{max} = 3.5 - 2.3 = \qty{1.2}{eV}$.}
  \trap{use one unit for every term: all joules or all eV. $\phi$ is lower-case phi;
    capital $\Phi$ on the sheet is magnetic flux.}
\end{formulacard}

\begin{formulacard}[notonsheet]{Threshold frequency and threshold wavelength}
  \cardformula{\phi = hf_0 \qquad \lambda_0 = \frac{hc}{\phi}}
  \begin{symbols}
    \sym{$f_0$}{threshold frequency: lowest frequency that releases electrons}{\si{Hz}}
    \sym{$\lambda_0$}{threshold wavelength: longest wavelength that releases electrons}{\si{m}}
    \sym{$\phi$}{work function of the metal}{\si{J}}
    \sym{$h$}{Planck's constant, \qty{6.626e-34}{J.s}}{\si{J.s}}
    \sym{$c$}{speed of light in a vacuum, \qty{3.00e8}{m.s^{-1}}}{\si{m.s^{-1}}}
  \end{symbols}
  \inwords{at the threshold, one photon has just enough energy to free an electron.}
  \whentouse{to decide whether light of a given frequency or wavelength releases
    electrons.}
  \whennottouse{to find how fast the electrons move. Use $K_\mathrm{max} = hf - \phi$.}
  \rearranged{$f_0 = \dfrac{\phi}{h}$ \qquad $\phi = \dfrac{hc}{\lambda_0}$
    \qquad $K_\mathrm{max} = h(f - f_0)$}
  \tensecond{$\phi = \qty{2.0}{eV} = \qty{3.204e-19}{J}$, so
    $f_0 = \dfrac{\num{3.204e-19}}{\num{6.626e-34}} = \qty{4.835e14}{Hz}$ and
    $\lambda_0 = \dfrac{c}{f_0} = \qty{620}{nm}$.}
  \derive{at the threshold the fastest electron has nothing left over:
    \begin{align*}
      0 &= hf_0 - \phi && \text{put $K_\mathrm{max} = 0$, $f = f_0$ in $K_\mathrm{max} = hf - \phi$}\\
      \phi &= hf_0 && \text{add $\phi$ to both sides}\\
      \lambda_0 &= \frac{c}{f_0} = \frac{hc}{\phi} && \text{use $v = f\lambda$ with $v = c$, and $f_0 = \dfrac{\phi}{h}$}
    \end{align*}}
  \trap{electrons are released only when $f \ge f_0$, which means
    $\lambda \le \lambda_0$. A wavelength \emph{longer} than $\lambda_0$ never
    works, however bright the light.}
\end{formulacard}

How many photons does a lamp or laser give out each second? A source of power $P$
gives out $P$ joules of light energy each second. Each photon carries $hf$. So divide
the energy each second by the energy of one photon.

\begin{formulacard}[notonsheet]{Number of photons each second}
  \cardformula{\frac{N}{t} = \frac{P}{hf} = \frac{P\lambda}{hc}}
  \begin{symbols}
    \sym{$\dfrac{N}{t}$}{number of photons given out each second}{\si{s^{-1}}}
    \sym{$P$}{power of the light source (light energy each second)}{\si{W}}
    \sym{$h$}{Planck's constant, \qty{6.626e-34}{J.s}}{\si{J.s}}
    \sym{$f$}{frequency of the light}{\si{Hz}}
    \sym{$\lambda$}{wavelength of the light}{\si{m}}
    \sym{$c$}{speed of light in a vacuum, \qty{3.00e8}{m.s^{-1}}}{\si{m.s^{-1}}}
  \end{symbols}
  \inwords{the light energy given out each second, divided by the energy of one
    photon.}
  \whentouse{a lamp or laser of known power that gives one colour (one frequency).}
  \whennottouse{light with many colours, or to find the energy of one photon (use
    $E = hf$).}
  \rearranged{$P = \dfrac{N}{t}\,hf$ \qquad $P = \dfrac{N}{t}\,\dfrac{hc}{\lambda}$}
  \tensecond{a \qty{5.0}{mW} green laser, $\lambda = \qty{532}{nm}$:
    $\dfrac{N}{t} = \dfrac{\num{5.0e-3} \times \num{5.32e-7}}{\num{6.626e-34} \times
    \num{3.00e8}} = \qty{1.3e16}{s^{-1}}$.}
  \derive{power is energy each second, $P = \dfrac{\Delta E}{\Delta t}$ (on the sheet).
    \begin{align*}
      P &= \frac{N}{t} \times hf && \text{(photons each second) $\times$ (energy of one photon)}\\
      \frac{N}{t} &= \frac{P}{hf} && \text{divide both sides by $hf$}\\
      \frac{N}{t} &= \frac{P\lambda}{hc} && \text{replace $f$ by $\dfrac{c}{\lambda}$}
    \end{align*}}
  \trap{change \si{mW} to \si{W} (divide by 1000) and \si{nm} to \si{m}. At the same
    power, red light gives \emph{more} photons each second than blue light, because
    each red photon has less energy.}
\end{formulacard}

\subsection{Planck, Einstein and a new model of light}

\Cref{ch3:tab:planckeinstein} compares Planck's and Einstein's quantum ideas.

\begin{table}
  \caption{Planck's and Einstein's quantum ideas.}\label{ch3:tab:planckeinstein}
  \centering
  \small
  \begin{tabular}{@{}L{32mm}L{52mm}L{52mm}@{}}
    \toprule
    & \colhead{Planck (1900)} & \colhead{Einstein (1905)}\\
    \midrule
    What is quantised? & energy of the oscillators in a black body's walls
      & light itself: it travels as photons\\
    Were quanta real? & no, he saw them as a maths trick & yes, a real property of light\\
    What did it explain? & the black body curve (no ultraviolet catastrophe)
      & the photoelectric effect\\
    The same idea & \multicolumn{2}{L{104mm}@{}}{energy is given out or taken in only
      in packets of $E = hf$}\\
    \bottomrule
  \end{tabular}
\end{table}

Light still shows interference and diffraction, which only waves explain
(\cref{ch2:chap}). But it is given out and absorbed as photons. So light has
\term{wave--particle duality}: it travels like a wave and gives up its energy like a
stream of particles. The photon is \emph{not} Newton's corpuscle model, which wrongly
predicted that light moves faster in water. Millikan's experiments (1914--1916)
confirmed Einstein's equation, and Einstein won the 1921 Nobel Prize for it.

\subsection{Photoelectric graphs}

\begin{toolkit}{Straight-line graphs}
\begin{itemize}
  \item A straight line is $y = mx + c$. Here $m$ is the gradient (how steep the line
        is) and $c$ is the $y$-intercept (where the line crosses the $y$-axis).
        Careful: here $c$ is \emph{not} the speed of light.
  \item A \term{line of best fit} is one straight line, drawn with a ruler. The
        points are spread evenly on both sides of it. It does not have to pass
        through any point.
  \item To \term{extrapolate} means to extend the line past your data. Draw this
        part dotted. Then you can read a value you did not measure, such as an
        intercept.
  \item Gradient $= \dfrac{\text{rise}}{\text{run}} = \dfrac{y_2 - y_1}{x_2 - x_1}$,
        using two points \emph{on the line of best fit}, far apart. Its unit is
        ($y$ unit)/($x$ unit). Read the power of ten from the axis label.
  \item The $x$-intercept is where $y = 0$. The $y$-intercept is where $x = 0$.
\end{itemize}
\end{toolkit}

Compare Einstein's equation with a straight line, $y = mx + c$:
\[
  \underbrace{K_\mathrm{max}}_{y} = \underbrace{h}_{m}\,\underbrace{f}_{x}
  \;\underbrace{-\,\phi}_{\text{$y$-intercept}}
\]
\Cref{ch3:fig:kf} shows the graph for two metals.

\begin{figure}
  \centering
  \begin{tikzpicture}
    \begin{axis}[hsc, width=0.74\linewidth, height=0.44\linewidth,
        axis lines=middle, xmin=0, xmax=14.5, ymin=-5, ymax=4.2,
        xtick={2,4,...,14}, ytick={-4,-2,2,4},
        xlabel={$f$ ($\times10^{14}$ Hz)}, ylabel={$K_\mathrm{max}$ (eV)},
        x label style={anchor=west}, y label style={anchor=south, rotate=-90},
        every axis x label/.style={at={(axis cs:14.5,0)}, anchor=west},
        every axis y label/.style={at={(axis cs:0,4.2)}, anchor=south},
        xticklabel style={font=\footnotesize, fill=white, inner sep=1pt},
        yticklabel style={font=\footnotesize}]
      \addplot[hscseriesA, line width=1.4pt, domain=5.077:14, samples=2] {0.4136*x - 2.10};
      \addplot[hscseriesA, line width=1pt, densely dotted, domain=0:5.077, samples=2] {0.4136*x - 2.10};
      \addplot[hscseriesB, line width=1.4pt, domain=10.40:14, samples=2] {0.4136*x - 4.30};
      \addplot[hscseriesB, line width=1pt, densely dotted, domain=0:10.40, samples=2] {0.4136*x - 4.30};
      \fill[hscseriesA] (axis cs:5.077,0) circle (2.2pt);
      \fill[hscseriesB] (axis cs:10.40,0) circle (2.2pt);
      \fill[hscseriesA] (axis cs:0,-2.10) circle (2.2pt);
      \fill[hscseriesB] (axis cs:0,-4.30) circle (2.2pt);
      \node[anchor=south east, font=\footnotesize, text=hscseriesA] at (axis cs:5.0,0.1) {$f_0$};
      \node[anchor=south east, font=\footnotesize, text=hscseriesB] at (axis cs:10.3,0.1) {$f_0$};
      \node[anchor=north west, font=\footnotesize, text=hscseriesA, fill=white, inner sep=1pt] at (axis cs:0.2,-2.2) {$-\phi = \qty{-2.10}{eV}$};
      \node[anchor=north west, font=\footnotesize, text=hscseriesB] at (axis cs:0.2,-4.4) {$-\phi = \qty{-4.30}{eV}$};
      \node[anchor=west, font=\small] at (axis cs:14.1,3.69) {caesium};
      \node[anchor=west, font=\small] at (axis cs:14.1,1.49) {zinc};
      \node[anchor=west, font=\footnotesize, align=left] at (axis cs:6.4,-3.4)
        {same gradient $h$: parallel lines.\\dotted: extrapolated (no electrons)};
    \end{axis}
  \end{tikzpicture}
  \caption{$K_\mathrm{max}$ against frequency for caesium and zinc.}\label{ch3:fig:kf}
\end{figure}

\begin{itemize}
  \item \textbf{Gradient} $= h$ for every metal, so the lines are parallel. With
        $K_\mathrm{max}$ in eV it is $h \div \num{1.602e-19} = \qty{4.136e-15}{eV.s}$.
  \item \textbf{$f$-intercept} $= f_0$. Below it there are no electrons.
  \item \textbf{$K_\mathrm{max}$-intercept} $= -\phi$. Extrapolate to $f = 0$ and
        change the sign. A bigger $\phi$ moves the line right and down.
\end{itemize}

Experiments measure $V_0$, not $K_\mathrm{max}$. So we turn Einstein's equation into
an equation for $V_0$. A graph of $V_0$ against $f$ is then also a straight line.

\begin{formulacard}[notonsheet]{Stopping voltage against frequency}
  \cardformula{V_0 = \frac{h}{q}f - \frac{\phi}{q}}
  \begin{symbols}
    \sym{$V_0$}{stopping voltage}{\si{V}}
    \sym{$h$}{Planck's constant, \qty{6.626e-34}{J.s}}{\si{J.s}}
    \sym{$q$}{size of the electron's charge, \qty{1.602e-19}{C}}{\si{C}}
    \sym{$f$}{frequency of the light}{\si{Hz}}
    \sym{$\phi$}{work function of the metal}{\si{J}}
  \end{symbols}
  \inwords{$V_0$ rises in a straight line with $f$. The gradient is $\dfrac{h}{q}$ for
    every metal.}
  \whentouse{to read $h$, $f_0$ and $\phi$ from a graph of $V_0$ against $f$.}
  \whennottouse{for a graph of $K_\mathrm{max}$ against $f$ (there the gradient is
    $h$ itself).}
  \rearranged{$h = q \times \text{gradient}$ \qquad $f_0 = f\text{-intercept}$
    \qquad $\phi = -q \times (V_0\text{-intercept})$}
  \tensecond{gradient \qty{4.08e-15}{V.s} and $V_0$-intercept \qty{-1.95}{V}:
    $h = \num{1.602e-19} \times \num{4.08e-15} = \qty{6.54e-34}{J.s}$, and
    $\phi = \qty{1.95}{eV}$.}
  \derive{start from Einstein's equation:
    \begin{align*}
      K_\mathrm{max} &= hf - \phi && \text{Einstein's equation (on the sheet)}\\
      qV_0 &= hf - \phi && \text{replace $K_\mathrm{max}$ by $qV_0$}\\
      V_0 &= \frac{h}{q}f - \frac{\phi}{q} && \text{divide every term by $q$}\\
      \frac{h}{q} &= \frac{\qty{6.626e-34}{J.s}}{\qty{1.602e-19}{C}}
        = \qty{4.136e-15}{V.s} && \text{the gradient for every metal}
    \end{align*}}
  \trap{the $V_0$-intercept is $-\dfrac{\phi}{q}$, in volts. Change the sign to get
    $\phi$ in eV. Multiply by $q$ to get $\phi$ in joules.}
\end{formulacard}

\Cref{ch3:fig:iv} shows photocurrent against the voltage across the tube.

\begin{sidebyside}[0.5]
  \centering
  \begin{tikzpicture}
    \begin{axis}[hsc, width=\linewidth, height=0.8\linewidth,
        axis lines=middle, xmin=-3, xmax=1.8, ymin=0, ymax=3.4,
        xtick={-2.4,-1.2}, xticklabels={$-V_{0}$ (A), $-V_{0}$ (B)}, ytick=\empty,
        xlabel={$V$}, ylabel={$I$ (photocurrent)},
        every axis x label/.style={at={(axis cs:1.8,0)}, anchor=west},
        every axis y label/.style={at={(axis cs:0,3.4)}, anchor=south},
        xticklabel style={font=\footnotesize}]
      \addplot[hscseriesA, line width=1.3pt, domain=-2.4:0, samples=40] {2.2*sin(90*(x+2.4)/2.4)^2};
      \addplot[hscseriesA, line width=1.3pt, domain=0:1.6, samples=2] {2.2};
      \addplot[hscseriesB, line width=1.3pt, dashed, domain=-1.2:0, samples=40] {1.1*sin(90*(x+1.2)/1.2)^2};
      \addplot[hscseriesB, line width=1.3pt, dashed, domain=0:1.6, samples=2] {1.1};
      \node[anchor=south, font=\footnotesize] at (axis cs:1.45,2.2) {A};
      \node[anchor=south, font=\footnotesize] at (axis cs:1.45,1.1) {B};
      \node[anchor=north, font=\scriptsize, align=center] at (axis cs:0.6,1.98)
        {saturation\\current (A)};
      \draw[->, hscink, line width=0.5pt] (axis cs:0.6,1.98) -- (axis cs:0.6,2.15);
      \node[anchor=south, font=\scriptsize, align=center, text=hscmuted]
        at (axis cs:-1.5,2.55) {anode negative\\(pushes electrons back)};
      \node[anchor=south, font=\scriptsize, align=center, text=hscmuted]
        at (axis cs:0.85,2.55) {anode\\positive};
    \end{axis}
  \end{tikzpicture}
  \captionof{figure}{Photocurrent against tube voltage (idealised) for two beams.}
  \label{ch3:fig:iv}
\nextside
  \begin{itemize}
    \item Positive voltage: every released electron reaches the anode. The flat
          part is the \term{saturation current}, set by the number of photons each
          second.
    \item Negative voltage: only faster electrons get through. $I = 0$ at $-V_0$.
    \item Beam A has more photons each second (higher flat part) and a higher
          frequency (bigger $V_0$) than beam B.
  \end{itemize}
\end{sidebyside}

Brighter light of the same $f$ raises only the flat part. Higher $f$ with the same
number of photons each second makes $V_0$ bigger. (At the same \emph{power},
higher-$f$ light has fewer photons, so a lower flat part.) A metal with a bigger
$\phi$, in the same light, gives a smaller $V_0$: its curve meets the voltage axis
closer to zero.

\begin{workedexample}{Will it work? Zinc and two lamps}\label{ch3:we:zinc}
Zinc ($\phi = \qty{4.30}{eV}$) is lit by a very bright \qty{365}{nm} lamp and a dim
\qty{254}{nm} lamp. (a) Find the threshold wavelength of zinc. (b) Which lamp releases
electrons? Explain. (c) For that lamp, find $K_\mathrm{max}$ and $V_0$.
\given{$\phi = 4.30 \times \num{1.602e-19} = \qty{6.889e-19}{J}$ \ (eV to J), \
  $\lambda_1 = \qty{3.65e-7}{m}$, \ $\lambda_2 = \qty{2.54e-7}{m}$}
\find{$\lambda_0$; which lamp works; $K_\mathrm{max}$ and $V_0$}
\formula{$\lambda_0 = \dfrac{hc}{\phi}$, \ $E = \dfrac{hc}{\lambda}$, \
  $K_\mathrm{max} = hf - \phi$, \ $K_\mathrm{max} = qV_0$}
\substitute{\eqlines{\lambda_0 &= \frac{\qty{6.626e-34}{J.s} \times \qty{3.00e8}{m.s^{-1}}}{\qty{6.889e-19}{J}}
    = \qty{2.886e-7}{m}}\par
  \eqlines{E_{365} &= \frac{\qty{6.626e-34}{J.s} \times \qty{3.00e8}{m.s^{-1}}}{\qty{3.65e-7}{m}}
    = \qty{5.446e-19}{J}\\
  &= \frac{\num{5.446e-19}}{\num{1.602e-19}} = \qty{3.40}{eV} \quad\text{(J to eV)}}\par
  \eqlines{E_{254} &= \frac{\qty{6.626e-34}{J.s} \times \qty{3.00e8}{m.s^{-1}}}{\qty{2.54e-7}{m}}
    = \qty{7.826e-19}{J}\\
  &= \frac{\num{7.826e-19}}{\num{1.602e-19}} = \qty{4.885}{eV} \quad\text{(J to eV)}}\par
  \eqlines{K_\mathrm{max} &= \qty{4.885}{eV} - \qty{4.30}{eV} = \qty{0.585}{eV}
    \quad\text{(all in eV)}\\
  V_0 &= \frac{K_\mathrm{max}}{q} = \frac{0.585 \times \qty{1.602e-19}{J}}{\qty{1.602e-19}{C}}
    = \qty{0.585}{V}}}
\answer{(a) \final{$\lambda_0 = \qty{2.89e-7}{m}$} (\qty{289}{nm}). (b) Only the dim
  \qty{254}{nm} lamp, because \qty{254}{nm} $< \qty{289}{nm}$. A \qty{365}{nm} photon
  has only \qty{3.40}{eV} $< \phi$, and a brighter lamp only gives more of them.
  (c) \final{$K_\mathrm{max} = \qty{0.59}{eV}$}, so \final{$V_0 = \qty{0.59}{V}$}.}
\checkstep{\qty{254}{nm} is shorter than $\lambda_0 = \qty{289}{nm}$ and
  \qty{365}{nm} is longer \cmark. \qty{4.89}{eV} is only a little above $\phi$, so a
  small $K_\mathrm{max}$ makes sense \cmark.}
\end{workedexample}

\begin{workedexample}{Finding \texorpdfstring{$h$}{h} from data (HSC style)}\label{ch3:we:data}
A student measures the stopping voltage for six frequencies of light on sodium.
\begin{center}\small
  \begin{tabular}{@{}lcccccc@{}}
    \toprule
    $f$ ($\times10^{14}$ Hz) & 6.0 & 7.0 & 8.0 & 9.0 & 10.0 & 11.0\\
    $V_0$ (V) & 0.21 & 0.58 & 1.02 & 1.40 & 1.85 & 2.24\\
    \bottomrule
  \end{tabular}
\end{center}
(a) Plot $V_0$ against $f$ with a line of best fit. (b) Use the gradient to find
Planck's constant. (c) Find the threshold frequency and the work function of sodium.
\given{the data; \ $q = \qty{1.602e-19}{C}$}
\find{$h$, $f_0$ and $\phi$}
\formula{$V_0 = \dfrac{h}{q}f - \dfrac{\phi}{q}$: \ $h = q \times \text{gradient}$,
  \ $f_0 = f$-intercept, \ $\phi = -q \times (V_0\text{-intercept})$}
\substitute{(a) See \cref{ch3:fig:data}. Two points far apart \emph{on the line}:
  ($\qty{6.0e14}{Hz}$, \qty{0.19}{V}) and ($\qty{11.0e14}{Hz}$, \qty{2.24}{V}).\par
  \eqlines{\text{gradient} &= \frac{\qty{2.24}{V} - \qty{0.19}{V}}{\qty{11.0e14}{Hz} - \qty{6.0e14}{Hz}}
    = \qty{4.10e-15}{V.s}\\
  h &= \qty{1.602e-19}{C} \times \qty{4.10e-15}{V.s} = \qty{6.568e-34}{J.s}}\par
  \emph{Intercept:} rearrange $y = mx + c$ to $c = y - mx$. Use the point
  (\qty{6.0e14}{Hz}, \qty{0.19}{V}) on the line.\par
  \eqlines{V_0\text{-intercept} &= \qty{0.19}{V} - \qty{4.10e-15}{V.s} \times \qty{6.0e14}{Hz}
    = \qty{-2.27}{V}}\par
  \emph{Threshold:} at the $f$-intercept, $V_0 = 0$. Add 2.27 to both sides, then
  divide by the gradient.\par
  \eqlines{0 &= (\qty{4.10e-15}{V.s})\,f_0 - \qty{2.27}{V}\\
  f_0 &= \frac{\qty{2.27}{V}}{\qty{4.10e-15}{V.s}} = \qty{5.537e14}{Hz}}\par
  \emph{Work function:} $\phi = -q \times (V_0\text{-intercept})$. Then change J to eV.
  It is the same number as the intercept, without the minus sign.\par
  \eqlines{\phi &= -\qty{1.602e-19}{C} \times (\qty{-2.27}{V}) = \qty{3.637e-19}{J}\\
  \phi &= \frac{\num{3.637e-19}}{\num{1.602e-19}} = \qty{2.27}{eV}}}
\begin{figure}
  \centering
  \begin{tikzpicture}
    \begin{axis}[hsc, width=0.62\linewidth, height=0.33\linewidth,
        axis lines=middle, xmin=0, xmax=12.5, ymin=-2.6, ymax=2.8,
        xtick={2,4,...,12}, ytick={-2,-1,1,2},
        xlabel={$f$ ($\times10^{14}$ Hz)}, ylabel={$V_0$ (V)},
        every axis x label/.style={at={(axis cs:12.5,0)}, anchor=west},
        every axis y label/.style={at={(axis cs:0,2.8)}, anchor=south},
        xticklabel style={font=\footnotesize, fill=white, inner sep=1pt},
        yticklabel style={font=\footnotesize}]
      \addplot[hscseriesB, line width=1.1pt, domain=5.537:12, samples=2] {0.41*x - 2.27};
      \addplot[hscseriesB, line width=1pt, densely dotted, domain=0:5.537, samples=2] {0.41*x - 2.27};
      \addplot[hscseriesA, only marks, mark=*, mark size=2pt]
        coordinates {(6.0,0.21) (7.0,0.58) (8.0,1.02) (9.0,1.40) (10.0,1.85) (11.0,2.24)};
      \draw[guide] (axis cs:6.0,0.19) -- (axis cs:11.0,0.19) -- (axis cs:11.0,2.24);
      \node[anchor=south, font=\footnotesize] at (axis cs:9.3,0.19) {run};
      \node[anchor=west, font=\footnotesize] at (axis cs:11.05,1.2) {rise};
      \fill[hscseriesB] (axis cs:5.537,0) circle (2pt);
      \node[anchor=south east, font=\footnotesize] at (axis cs:5.3,0.15) {$f_0 \approx \qty{5.5e14}{Hz}$};
      \fill[hscseriesB] (axis cs:0,-2.27) circle (2pt);
      \node[anchor=north west, font=\footnotesize] at (axis cs:0.2,-2.3) {\qty{-2.27}{V}};
    \end{axis}
  \end{tikzpicture}
  \caption{Stopping voltage against frequency for sodium (dotted: extrapolated).}
  \label{ch3:fig:data}
\end{figure}
\answer{(b) \final{$h = \qty{6.57e-34}{J.s}$}. (c) \final{$f_0 = \qty{5.54e14}{Hz}$}
  and \final{$\phi = \qty{2.27}{eV} = \qty{3.64e-19}{J}$}.}
\checkstep{$h$ is 0.9\,\% below \qty{6.626e-34}{J.s} \cmark. Sodium's work function
  is about \qty{2.3}{eV} \cmark.}
\end{workedexample}

\begin{mistakes}
\begin{itemize}
  \item \xmark\ Subtracting an energy in eV from an energy in J (for example $hf$ in J
        minus $\phi$ in eV). \cmark\ Change both to the same unit first.
  \item \xmark\ ``Bright red light will release electrons from any metal.''
        \cmark\ If $hf < \phi$, none leave, however bright the light.
  \item \xmark\ ``The wavelength must be longer than $\lambda_0$.''
        \cmark\ It must be \emph{shorter} (or equal): $\lambda \le \lambda_0$.
  \item \xmark\ Putting $K_\mathrm{max}$ in eV into $K = \frac12 mv^2$.
        \cmark\ Change it to joules first.
  \item \xmark\ ``The intercept is $\phi$'', or ``each metal has its own gradient''.
        \cmark\ The intercept is $-\phi$ (change the sign). Every metal has the
        same gradient.
\end{itemize}
\end{mistakes}

\begin{examtip}
The 2025 HSC exam gave a table of stopping voltage against frequency for one metal
(6 marks).
\begin{enumerate}
  \item Graph and threshold frequency (3 marks): if the axes are not labelled,
        label them with units. Plot every point accurately. Draw one straight line of best fit (not dot-to-dot). Extrapolate it
        (dotted) to the frequency axis. Read $f_0$ with its unit.
  \item ``Using the particle model, explain the features'' (3 marks): light is a
        stream of photons, each with energy $E = hf$. One photon gives all its energy
        to one electron. By conservation of energy, $qV_0 = K_\mathrm{max} = hf -
        \phi$, so $V_0$ rises in a straight line with $f$ (gradient $\dfrac{h}{q}$).
        At the threshold, $hf_0 = \phi$. Below $f_0$ no single photon has enough
        energy to free an electron, so the current is zero, however bright the light.
\end{enumerate}
\end{examtip}

\begin{quickcheck}
  \item Sodium has a work function of \qty{2.28}{eV}. Calculate its threshold
        wavelength. Will yellow light of wavelength \qty{589}{nm} release electrons?
  \item Ultraviolet light of wavelength \qty{200}{nm} falls on copper
        ($\phi = \qty{4.70}{eV}$). Find $K_\mathrm{max}$ (in eV and J), the stopping
        voltage and the maximum speed of the photoelectrons.\label{ch3:qc:copper}
  \item A \qty{2.0}{mW} violet laser (\qty{405}{nm}) shines on caesium
        ($\phi = \qty{2.10}{eV}$). (a) How many photons arrive each second? (b) The
        power is doubled, at the same wavelength. What happens to the number of
        photoelectrons each second and to $K_\mathrm{max}$? Describe how the
        current--voltage graph changes. Explain.
  \item Metals P and Q give parallel lines on a $K_\mathrm{max}$--$f$ graph, crossing
        the frequency axis at \qty{6.0e14}{Hz} (P) and \qty{9.0e14}{Hz} (Q). Which
        has the larger work function? Calculate the work function of P in eV.
  \item Compare Planck's and Einstein's ideas about quanta. Give one similarity and
        two differences. (3 marks)\label{ch3:qc:planckeinstein}
\end{quickcheck}

%% ==================================================================
\section{Key experiments}\label{ch3:sec:experiments}
\covers{M7.IQ3.01, M7.IQ3.02, M7.IQ3.03}

The syllabus asks you to analyse the black body evidence and to investigate the
photoelectric evidence. For each experiment, learn its aim, method and results, and
how to judge its validity, reliability and accuracy.

\begin{experiment}{Measuring black body radiation (1899--1900)}
\aim{To measure how the intensity of black body radiation depends on wavelength, at
  several temperatures (Lummer and Pringsheim, Berlin).}
\apparatus{an electric oven with a hollow inside and a small hole (the black body); a
  thermocouple (thermometer); a slit; a prism to spread the radiation into
  wavelengths; a \term{bolometer}, a detector whose resistance changes when
  radiation warms it (\cref{ch3:fig:bbapp}).}
\method{\begin{enumerate}
  \item Heat the oven to a steady temperature, measured with the thermocouple.
  \item Pass the radiation from the hole through the slit and the prism.
  \item Move the bolometer along the spectrum and record the intensity at each
        wavelength. Repeat at several temperatures.
\end{enumerate}
  \begin{figure}
    \centering
    \begin{tikzpicture}[x=1cm, y=1cm]
      \fill[hscmuted!40] (0,0) rectangle (2.3,1.7);
      \fill[hscink!85] (0.35,0.35) rectangle (1.95,1.35);
      \fill[white] (1.94,0.75) rectangle (2.31,0.95);
      \draw[hscforce, line width=0.9pt] (0.9,1.9) -- (0.9,1.0);
      \node[diagramlabel, anchor=south, font=\footnotesize] at (0.9,1.9) {thermocouple};
      \node[diagramlabel, anchor=north, font=\footnotesize] at (1.15,-0.05)
        {hot oven (black body)};
      \draw[hscink, line width=2pt] (3.0,0.35) -- (3.0,0.78) (3.0,0.92) -- (3.0,1.35);
      \node[diagramlabel, anchor=north, font=\footnotesize] at (3.0,0.3) {slit};
      \draw[ray] (2.3,0.85) -- (4.1,0.85);
      \draw[hscink, line width=0.8pt, fill=hscbrand!10] (3.8,0.25) -- (5.0,0.25) -- (4.4,1.45) -- cycle;
      \draw[hscray, line width=1pt] (4.1,0.85) -- (4.73,0.79);
      \node[diagramlabel, anchor=north, font=\footnotesize] at (4.4,0.2) {prism};
      \draw[hscseriesB, line width=1pt] (4.73,0.80) -- (7.4,0.55);
      \draw[hscseriesC, line width=1pt] (4.73,0.79) -- (7.4,0.15);
      \draw[hscvel, line width=1pt] (4.73,0.78) -- (7.4,-0.25);
      \node[diagramlabel, anchor=west, font=\footnotesize] at (7.45,0.55) {longer $\lambda$ (bent least)};
      \node[diagramlabel, anchor=west, font=\footnotesize] at (7.45,-0.25) {shorter $\lambda$ (bent most)};
      \draw[hscink, line width=0.8pt, fill=white] (6.35,0.21) rectangle (6.65,0.51);
      \draw[hscmuted, line width=0.5pt] (6.5,0.51) -- (6.2,1.9)
        node[diagramlabel, anchor=south, font=\footnotesize] {bolometer (moved along the spectrum)};
    \end{tikzpicture}
    \caption{Measuring a black body curve (simplified).}\label{ch3:fig:bbapp}
  \end{figure}}
\results{curves like \cref{ch3:fig:bbcurves}. Typical peaks (example readings made
  up for this book, not the original data): \qty{2.92}{\micro m} at \qty{1000}{K},
  \qty{2.30}{\micro m} at \qty{1250}{K}, \qty{1.94}{\micro m} at \qty{1500}{K}, so
  $\lambda_\mathrm{max}T \approx \qty{2.9e-3}{m.K}$ each time.}
\conclusion{$\lambda_\mathrm{max}T$ is constant: Wien's law. Classical theory
  fitted only the long-wavelength side; Planck's formula fitted the whole curve.}
\significance{these curves led Planck to propose quanta (1900), the first step to
  the photon model of light.}
\validity{\begin{itemize}
  \item \textbf{Validity:} a small hole, so the oven acts as a black body; a steady
        temperature measured inside the cavity; a prism that lets all the measured
        wavelengths through (glass absorbs much infrared).
  \item \textbf{Reliability:} repeat each scan at the same temperature; the curves
        should agree.
  \item \textbf{Accuracy:} calibrate the thermocouple and bolometer; use a narrow
        slit; shield the bolometer from room heat.
\end{itemize}}
\end{experiment}

\begin{experiment}{The photoelectric effect: measuring \texorpdfstring{$h$}{h} and \texorpdfstring{$\phi$}{phi}}
\aim{To find how $V_0$ depends on frequency and brightness, and to find $h$ and
  $\phi$ (Lenard 1902; Millikan 1914--1916; a school photocell or a simulation).}
\apparatus{a photocell in the circuit of \cref{ch3:fig:apparatus} (reversible supply,
  voltmeter, sensitive ammeter); light of known frequencies (a mercury lamp with
  colour filters, or LEDs). The free PhET ``Photoelectric Effect'' computer
  simulation can be used instead.}
\method{\begin{enumerate}
  \item Shine light of one frequency on the cathode. Make the anode negative.
  \item Increase the reverse voltage until the current just becomes zero. Record
        $V_0$. Repeat three times and average. Do this for at least five frequencies.
  \item Keep $f$ fixed. Move the lamp closer or further away to change the
        brightness. For each brightness, make the anode positive and record the
        saturation current. Then make the anode negative again and find $V_0$.
\end{enumerate}}
\results{$V_0$ rises in a straight line with $f$ (like \cref{ch3:fig:data}). No current
  below $f_0$, however bright the light. Brighter light raises the saturation current
  but not $V_0$.}
\conclusion{$qV_0 = hf - \phi$: $q \times$ gradient gives $h$; the intercepts give
  $f_0$ and $\phi$. Brightness changes only the number of electrons. This supports
  the photon model, not the wave model.}
\significance{Lenard's result, that $K_\mathrm{max}$ does not depend on brightness,
  was key evidence the wave model could not explain. Millikan expected to prove
  Einstein wrong. Instead he confirmed the straight-line law and found $h$ within
  about 1\,\% of today's value.}
\validity{\begin{itemize}
  \item \textbf{Validity:} change only the frequency (same metal and tube); use light
        of nearly one colour; block room light.
  \item \textbf{Reliability:} repeat each $V_0$ three times; use five or more
        frequencies; check the scatter about the line.
  \item \textbf{Accuracy:} the current falls to zero slowly, so $V_0$ is hard to judge
        (use a very sensitive ammeter). Keep the metal clean: an oxide layer (a thin
        coating that forms when the metal reacts with air) changes $\phi$. Use the
        gradient, not one single reading. An error that shifts every $V_0$ by the
        same amount does not change the gradient, so $h$ is still accurate.
\end{itemize}}
\end{experiment}

\begin{examtip}
Questions often ask you to ``assess the validity, reliability or accuracy'' of one of
these experiments. Answer with the right word and a reason:
\begin{itemize}
  \item \textbf{Validity:} name what was kept the same (same metal, same tube, one
        colour at a time), so only the variable being tested changed.
  \item \textbf{Reliability:} say that readings were repeated and agreed closely.
  \item \textbf{Accuracy:} compare your result with the accepted value, for example
        $h$ with \qty{6.626e-34}{J.s}, as a percentage difference.
\end{itemize}
\end{examtip}
